\scanpage{74}
\mspage{73}

C'est la donnée d'une application
\[
M\longrightarrow\unsure{\mathcal{P}|\mathscr{P}}(P)\qquad
X\unsure{\mathrel{\vert\!{-}\!\vert}|\longmapsto}|X|
\]
satisfaisant les conditions
\[
X\ll Y\quad\Longrightarrow\quad |X|\subset |Y|
\tag*{$(\mathrm{i})_M$}
\]
\margin{\uline{application croissante} pour $\ll$, $\subset$}
\[
X\unsure{\bowtie|\mathrel{\gtrless}}Y\quad\Longrightarrow\quad
|X|\cap|Y|=\bigcup_{Z\in\widetilde X\cap\widetilde Y}|Z|
\tag*{$(\mathrm{ii})_M$}
\]
\[
X\parallel Y\quad\Longleftrightarrow\quad |X|\cap|Y|=\varnothing
\tag*{\tikz[baseline=(n.base)]{\node[draw,circle,inner sep=1pt](n){$(\mathrm{iii})_M$};}}
\]
\margin{$\longleftarrow$\ \unsure{le fort} de l'axiome est dans l'implication $\Longleftarrow$}
\margin{(NB on a $\Longrightarrow$ par (ii))}
\[
\deleted{\unsure{X}}\quad X\unsure{\preceq|\leq}Y\unsure{,|\ast}\quad\Longrightarrow\quad |X|=|Y|
\tag*{$(\deleted{\unsure{\mathrm{iii}}}\mathrm{iv})_M$}
\]
\[
\unsure{\backslash}\qquad |X|=\bigcup_{Z\in\unsure{\operatorname{omb}|\operatorname{Somb}|\operatorname{Comb}}(X)}|Z|
\tag*{$(\mathrm{v})_M$}
\]
\[
\begin{aligned}
&\inserted{\text{\tikz[baseline=(n.base)]{\node[draw,circle,inner sep=0.7pt](n){a};}\ }\bigcup_{X\in M}|X|=P,\ \text{et}}\\
&\inserted{\text{\tikz[baseline=(n.base)]{\node[draw,circle,inner sep=0.7pt](n){b};}}}\ \text{Si }x,y\in P,\quad\exists X\in M\ \text{avec}\quad x\in|X|,\ y\notin|X|\\
&\hspace{26em}\text{ou l'inverse}
\end{aligned}
\tag*{$(\unsure{?\mathrm{vi}|\mathrm{vi}})_M$}
\]

NB 1) L'axiome (vi) est dans la nature d'une normalisation; il suffit de remplacer $P$ par son quotient par la relation d'équivalence
\[
x\sim y\quad\Longleftrightarrow\quad M_x=M_y
\qquad (\text{où }M_x=\{X\in M\mid x\in|X|\})
\]

\uline{NB} 2) La condition v implique qu'il suffit de connaître $|Z|$ pour $Z\in L$. \deleted{Mais cette \unsure{fonction}} On \uline{définit} les $|X|$ pour $X$ \unsure{général} par (v). Mais la validité des conditions (i) à (iv) \unsure{requiert|revient} des conditions draconiennes sur $|Z|$. La condition \unsure{(iii)|(ii)} p. ex. équivaut à la $=$ condition \inserted{(ou \unsure{(iii')|(ii')})} pour $X,Y\in L$, $+$ la condition At L2. La condition \unsure{(vi)|(iii)} \deleted{\unsure{doit être automatique} \unsure{moyennant} $L3$ \unsure{la condition} $(\unsure{\mathrm{ii}'|\mathrm{iii}'})$: $\unsure{Z\,X}$,} est automatique moyennant At L3 et \unsure{(ii')|(iii')} (partie $\Rightarrow$). La condition (i) tautologique. Je soupçonne que (iv) résulte des autres conditions. Donc, \unsure{ça revient} \unsure{via l'appl.|à l'appl.} \inserted{(moyennant At L2, L3)}
\[
L\longrightarrow\mathcal{P}(P),\qquad
x\mathrel{\vert\!{-}\!\vert}|x|,
\]
